
%  PLANTILLA GUÍAS INSTITUCIONALES — USACH PREMIUM
%  Fusiona el estilo formal de guia1 (portada + marca de agua)
%  con el estilo premium de guia3 (colores + tcolorbox + multicol).
%
%  INSTRUCCIONES RÁPIDAS:
%  1. Edita el bloque "INFORMACIÓN DE LA GUÍA" (líneas ~65-85).
%  2. Rellena \listacontenidos con los temas de tu guía.
%  3. Agrega ejercicios en los entornos indicados.
%  4. Compila con: pdflatex guia-institucional.tex (2 veces).
%
%  ESTRUCTURA DEL DOCUMENTO:
%    Portada → Introducción → Kit de Propiedades → Ejercicios → Solucionario
% ============================================================

\documentclass[letterpaper, 12pt]{article}

% ============================================================
% PAQUETES
% ============================================================
\usepackage[utf8]{inputenc}
\usepackage[spanish]{babel}
\usepackage{amsmath, amssymb}
\usepackage{geometry}
\usepackage{fancyhdr}
\usepackage{enumitem}
\usepackage{graphicx}
\usepackage{hyperref}
\usepackage{multicol}
\usepackage{parskip}
\usepackage{xcolor}
\usepackage{tcolorbox}
\usepackage{eso-pic}
\usepackage{transparent}
\usepackage{tikz}
\usetikzlibrary{patterns}
\tcbuselibrary{skins}

% ============================================================
% GEOMETRÍA
% ============================================================
\geometry{letterpaper, margin=1.5cm, top=3cm, bottom=2.5cm, headheight=2cm}

% ============================================================
% COLORES INSTITUCIONALES
% ============================================================
\definecolor{celesteSuave}{HTML}{F0F7FF}
\definecolor{azulFuerte}{HTML}{1A5276}
\definecolor{turquesaUsach}{HTML}{33CCCC}
\definecolor{enlace}{HTML}{1A5276}

% ============================================================
% RUTAS DE IMÁGENES  ← ajusta si tu carpeta de imágenes es distinta
% ============================================================
\newcommand{\logoup}{./images/logo_up.png}
\newcommand{\logousach}{./images/logo_usach.png}
\newcommand{\logofondo}{./images/logo_up.png}

% ============================================================
% INFORMACIÓN DE LA GUÍA  ← EDITA AQUÍ ANTES DE COMPILAR
% ============================================================
\newcommand{\institucion}{Usach Premium}
\newcommand{\curso}{Calculo I}              % ej: Cálculo I
\newcommand{\guiasnumero}{Guía N - Título}         % ej: Guía 1 - Trigonometría
\newcommand{\semestre}{1er. Semestre de 2026}
\newcommand{\preparadopor}{Usach Premium.}         % o el nombre del ayudante
\newcommand{\webinstitucion}{www.usachpremium.cl}
\newcommand{\instagraminstitucion}{https://www.instagram.com/usach.premium}
\newcommand{\transparencia}{0.05}                  % opacidad marca de agua (0.0–1.0)
\newcommand{\fraseportada}{"Por y para estudiantes"}

% Lista de contenidos que aparece en la portada dentro del tcolorbox.
% Agrega o quita \item según los temas de tu guía.
\newcommand{\listacontenidos}{
    \item[\color{azulFuerte}$\Rightarrow$] Dominio y rango de funciones
    \item[\color{azulFuerte}$\Rightarrow$] ¿Dónde crece y dónde decrece una función?
    \item[\color{azulFuerte}$\Rightarrow$]Composicion de funciones
    \item[\color{azulFuerte}$\Rightarrow$]Inyectividad,sobretectividad,Biyectividad
    \item[\color{azulFuerte}$\Rightarrow$] Funcion inversa
   
    
}

% ============================================================
% HYPERREF
% ============================================================
\hypersetup{
    colorlinks=true,
    linkcolor=enlace,
    urlcolor=enlace,
    citecolor=enlace,
    pdfborder={0 0 0},
}

% ============================================================
% ENCABEZADO Y PIE DE PÁGINA (páginas interiores)
% ============================================================
\pagestyle{fancy}
\fancyhf{}
\fancyhead[L]{%
    \includegraphics[height=1.2cm]{\logoup}%
}
\fancyhead[R]{%
    \begin{minipage}[b]{0.42\textwidth}
        \raggedleft
        \fontsize{9pt}{11pt}\selectfont
        \textbf{\color{azulFuerte}\institucion} \\
        \curso\ — \guiasnumero \\
        \textit{\color{gray}@usach.premium}
    \end{minipage}\hspace{0.1cm}%
    \includegraphics[height=1.2cm]{\logousach}%
}
\fancyfoot[L]{\fontsize{9pt}{11pt}\selectfont \textit{\fraseportada}}
\fancyfoot[R]{\fontsize{9pt}{11pt}\selectfont Página \thepage}
\renewcommand{\headrulewidth}{0.4pt}
\renewcommand{\footrulewidth}{0.4pt}

% Estilo portada: sin encabezado ni pie
\fancypagestyle{plain}{
    \fancyhf{}
    \renewcommand{\headrulewidth}{0pt}
    \renewcommand{\footrulewidth}{0pt}
}

% ============================================================
% MARCA DE AGUA (logo transparente al centro de cada página)
% Llama a \MarcaAgua dentro de un entorno para activarla.
% ============================================================
\newcommand{\MarcaAgua}{
    \AddToShipoutPicture{
        \AtPageCenter{
            \makebox(0,0){%
                \transparent{\transparencia}%
                \includegraphics[width=0.7\textwidth]{\logofondo}%
            }
        }
    }
}

% ============================================================
% CAJAS REUTILIZABLES
% ============================================================

% Caja de resumen / kit de propiedades
% Uso: \begin{kitpropiedades}{Título de la caja} ... \end{kitpropiedades}
\newenvironment{kitpropiedades}[1]{%
    \begin{tcolorbox}[colback=celesteSuave, colframe=azulFuerte,
        title=\textbf{#1}, fonttitle=\bfseries]
}{%
    \end{tcolorbox}
}

% Caja de desafío con sombra y título de color blanco sobre azul
% Uso: \begin{desafiobox}{Título} ... \end{desafiobox}
\newtcolorbox{desafiobox}[1]{
    colback=white,
    colframe=azulFuerte,
    fonttitle=\bfseries,
    coltitle=white,
    title=#1,
    arc=8pt,
    boxrule=1.2pt,
    enhanced,l
    drop shadow={black!5}
}

% Caja de bloque de soluciones
% Uso: \begin{solbox}{Título} ... \end{solbox}
\newenvironment{solbox}[1]{%
    \begin{tcolorbox}[colback=celesteSuave, colframe=azulFuerte,
        title=\textbf{#1}, fonttitle=\bfseries]
}{%
    \end{tcolorbox}
}

% ============================================================
% ENTORNOS DE SECCIÓN
% Cada entorno activa la marca de agua en las páginas que lo contienen.
% ============================================================

\newenvironment{introduccion}{%
    \section*{Introducción}%
    \MarcaAgua
}{}

\newenvironment{ejercicios}{%
    \MarcaAgua
}{}

\newenvironment{soluciones}{%
    \newpage
    \begin{center}
        \begin{tcolorbox}[colback=azulFuerte, colframe=azulFuerte,
            colupper=white, width=0.8\textwidth, center, arc=10pt]
            \centering\Large\textbf{SOLUCIONARIO}
        \end{tcolorbox}
    \end{center}
    \MarcaAgua
    \small
}{}

% ============================================================
\begin{document}
\thispagestyle{plain}

% ============================================================
% PORTADA
% ============================================================

% Logos en la portada: UP arriba a la izquierda, USACH arriba a la derecha
\AddToShipoutPicture*{%
    \put(54, 700){\includegraphics[width=2cm]{\logoup}}
    \put(420, 706){\includegraphics[width=5cm]{\logousach}}
}

\vspace*{0.5cm}

\begin{center}
    \color{azulFuerte}
    {\huge \textbf{GUÍA DE EJERCICIOS}} \\[0.5cm]
    {\Huge \textbf{\curso}} \\[0.3cm]
    {\Large \guiasnumero} \\[1.5cm]

    % Caja de contenidos
    \begin{tcolorbox}[colback=celesteSuave, colframe=azulFuerte, arc=10pt,
        width=0.85\textwidth, center, boxrule=1.5pt]
        \centering
        \vspace{0.2cm}
        {\Large \textbf{Contenidos}}
        \vspace{0.3cm}
        \color{black}
        \begin{itemize}[leftmargin=3cm]
            \listacontenidos
        \end{itemize}
        \vspace{0.2cm}
    \end{tcolorbox}
\end{center}

\vspace{1.5cm}

\begin{center}
    {\Large \textbf{\semestre}} \\[0.8cm]
    {\large \textbf{\preparadopor}}
\end{center}

\vspace{2cm}

% Pie de portada con links
\begin{center}
    {\color{azulFuerte}\hrule height 0.5pt}
    \vspace{0.3cm}
    \begin{minipage}{0.45\textwidth}
        \centering
        \textbf{Instagram:} \\
        \small\href{\instagraminstitucion}{@usach.premium}
    \end{minipage}
    \begin{minipage}{0.45\textwidth}
        \centering
        \textbf{Web:} \\
        \small\href{http://\webinstitucion}{\webinstitucion}
    \end{minipage}
\end{center}

\pagenumbering{arabic}
\newpage

% ============================================================
% ============================================================
% EJERCICIOS
% ============================================================


% --- Sección con ejercicios cortos ---
\subsection*{I. Ejercicios Propuestos}

\begin{enumerate}[label=\color{azulFuerte}\bfseries 1.\arabic*., leftmargin=*, itemsep=1em]

\item \textbf{Halle el dominio de la siguiente función:}
  $$h(w)=\frac{1}{\sqrt{16-w^{2}}}$$

\item Hallar el dominio y recorrido de las siguientes funciones:
    \begin{align*}
        F(t) &= \frac{1}{\sqrt{t}} \\[1ex]
        g(z) &= \frac{1}{\sqrt{4-z^{2}}}
    \end{align*}

\item Determina los intervalos donde cada función es creciente y donde es decreciente:
    \begin{itemize}
              \item $g(x) = (x-3)^{2}+2$
              \item $g(x) = -(x+1)^{2}+5$
    \end{itemize}

\item Considere la función $g(x)=\sqrt{3-\sqrt{x+2}}$:
    \begin{itemize}
        \item Encuentre el dominio natural de $g$.
        \item Demostrar que $g$ es estrictamente decreciente en todo su dominio.
        \item Encuentra el $\operatorname{Rec}(g)$.
    \end{itemize}

\end{enumerate}

\vspace{1cm}
\newpage

% --- Sección con ejercicios medianos (2 columnas) ---
\subsection*{II.Ejercicios propuestos}

2.2)
Considere las funciones $f(x)=\sqrt{x+1}$ y $g(x)=\sqrt{x-1}$. Para las siguientes funciones dadas, realice lo siguiente:
    \begin{itemize}
        \item Encuentre el dominio de la función resultante $(f\circ g)(x)$.
        \item Calcule su imagen en $5$ vía $(f\circ g)(x)$.
        \item Calcule su preimagen de $3$ vía $(f\circ g)(x)$ si es que existe.
    \end{itemize}
    \medskip
2.3)Considere las funciones $f(x)=\sqrt{\dfrac{x+1}{x-3}}$ y $g(x)=2x-1$:
    \begin{itemize}
        \item Determine el dominio de la composición $(f\circ g)(x)$.
        \item Calcule la imagen de $4$ para la función compuesta $(f\circ g)(x)$.
        \item Determine la preimagen de $1$ para la función $(f\circ g)(x)$ y vea si existe.
    \end{itemize} 
\medskip
2.4) Determine el valor de $k$ de modo que $f\circ g = g\circ f$, donde:
    \[
    f(x) = \frac{2x+3}{4} \quad \text{y} \quad g(x) = \frac{x-k}{2}
    \]
\medskip
2.5 Determine el valor de la constante $b$ de modo que se cumpla la condición $f\circ g = g\circ f$, dadas las funciones:
    \[
    f(x) = \frac{4x-1}{3} \quad \text{y} \quad g(x) = \frac{x+b}{4}
    \]
\medskip
2.6Considere la función $g:[1,+\infty[\,\rightarrow\mathbb{R}$, definida por:
    \[
    g(x) = 3x^{2}-6x+10
    \]
    Demuestre que $g$ es inyectiva, pero no sobreyectiva.
\vspace{4 cm}

\newpage
2.7) Sea $f:]0,2[\,\rightarrow\mathbb{R}$ la función dada por:
    \[
    f(x) = \frac{1}{2-x}-\frac{1}{x}
    \]
    Estudie la biyectividad de $f$ y, en caso de serlo, calcule su función inversa.
\medskip

2.8)Sea la función $f:[2,+\infty[\,\rightarrow[-5,+\infty[$ definida por:
    \[
    f(x) = 2x^{2}-8x+3
    \]
    Demuestre que $f$:
    \begin{itemize}
        \item Es inyectiva.
        \item Demuestre que $f$ es sobreyectiva.
        \item Determine la expresión de su función inversa $f^{-1}(x)$.
    \end{itemize}
\newpage
\section*{III. PEPs Antiguas}

\subsection*{0.1. Pep 2023}

\begin{enumerate}
    \item[\textbf{1.0)}] Dadas las funciones $f(x)=\sqrt{x-2}$ y $g(x)=\sqrt{25-x^{2}}-2$:
    \begin{itemize}
        \item Encuentra el dominio de $(f\circ g)(x)$.
    \end{itemize}

    \item[\textbf{1.1)}] Sea $f:A\rightarrow B$ una función definida como $f(x)=-x^{2}+4x-3$.
    \begin{itemize}
        \item Determine los conjuntos $A$ y $B$ de modo que $f$ sea biyectiva y encuentre $f^{-1}$.
    \end{itemize}
\end{enumerate}

\subsection*{0.2. Pep 2022}

\begin{enumerate}
    \item[\textbf{1.3)}] Considere la función $f:A\subset\mathbb{R}\rightarrow B\subset\mathbb{R}$ dada por:
    \[
    f(x) = -\sqrt{2-\sqrt{-1-2x}}
    \]
    \begin{itemize}
        \item[\textbf{a)}] Encontrar el conjunto $A=\text{Dom}(f)$.
        \item[\textbf{b)}] Determine si $f$ es creciente o decreciente.
        \item[\textbf{c)}] Determine $\text{Rec}(f)$.
    \end{itemize}
\end{enumerate}

\subsection*{0.3. Pep 2023 (Alternativa)}

\begin{enumerate}
    \item[\textbf{1.4)}] Sea $f:A\rightarrow B$ una función definida como:
    \[
    f(x) = x^{2}+8x+14
    \]
    Determine los conjuntos $A$ y $B$ de modo que $f$ sea biyectiva y que $x=1$ sea un elemento del $\text{Dom}(f)$. Luego encuentre $f^{-1}$.

    \item[\textbf{1.5)}] Considere la función $f(x)=1+\sqrt{\dfrac{2x-3}{x-4}}$:
    \begin{itemize}
        \item[\textbf{a)}] Determine $\text{Dom}(f)$ y $\text{Rec}(f)$.
        \item[\textbf{b)}] Determine la preimagen de $3$.
    \end{itemize}
1.6) Considere $f:]-1,1[\,\rightarrow\mathbb{R}$ definida como:
    \[
    f(x) = \frac{x}{1-|x|}
    \]
    Probar que $f$ es biyectiva y determinar $f^{-1}$.
\end{enumerate}
\newpage
\subsection*{ Más PEP}

\begin{enumerate}
    \item[\textbf{[1.8]}] Sea $k$ un número real de modo que la función $f:[-1,1]\rightarrow[k,10]$ dada por:
    \[
    f(x) = (x-2)^{2}+1
    \]
    está bien definida.
    \begin{itemize}
        \item[\textbf{a)}] Demuestre que $f$ es inyectiva.
        \item[\textbf{b)}] Determine el parámetro $k$ de modo que $f$ sea sobreyectiva.
        \item[\textbf{c)}] Determine la función inversa de $f$, indicando expresamente dominio, codominio y regla de asignación.
    \end{itemize}

    \item[\textbf{1.8)}] Sea:
    \[
    f(x) = \sqrt{4-\sqrt{2-3x}}
    \]
    \begin{itemize}
        \item[\textbf{a)}] Determine $\text{Dom}(f)$.
        \item[\textbf{b)}] Determine si la función es creciente o decreciente en su dominio.
    \end{itemize}

    \item[\textbf{1.9)}] Considere la fórmula:
    \[
    f(x) = \sqrt{\frac{3}{\sqrt{x+1}}-1}
    \]
    \begin{itemize}
        \item[\textbf{a)}] Halle el dominio natural de $f(x)$.
        \item[\textbf{b)}] Halle el recorrido de $f$.
        \item[\textbf{c)}] Si $g:\mathbb{R}\rightarrow\mathbb{R}$ es la función lineal $g(x)=x-2$, determine la función compuesta $f\circ g$ indicando expresamente su dominio y regla de asignación.
    \end{itemize}
\end{enumerate}
\newpage
% Diagrama TikZ opcional:
% \begin{center}
%   \begin{tikzpicture}[scale=1.2, thick]
%     % ... tu diagrama aquí ...
%   \end{tikzpicture}
% \end{center}



% ============================================================
% SOLUCIONARIO
% ============================================================
\begin{soluciones}

% ── Sección I ─────────────────────────────────────────────────
\begin{solbox}{I. Soluciones — Dominio, Recorrido y Crecimiento/Decrecimiento}
\begin{enumerate}[label=\color{azulFuerte}\bfseries 1.\arabic*., leftmargin=*, itemsep=0.5em]

    \item $\operatorname{Dom}(h) = (-4,\,4)$

    \item $\operatorname{Dom}(F)=(0,+\infty),\quad \operatorname{Rec}(F)=(0,+\infty)$ \\[2pt]
          $\operatorname{Dom}(g)=(-2,\,2),\quad \operatorname{Rec}(g)=\bigl[\tfrac{1}{2},+\infty\bigr)$

    \item $g(x)=(x-3)^2+2$: decreciente en $(-\infty,3]$, creciente en $[3,+\infty)$ \\[2pt]
          $g(x)=-(x+1)^2+5$: creciente en $(-\infty,-1]$, decreciente en $[-1,+\infty)$

    \item $\operatorname{Dom}(g)=[-2,\,7]$;\quad $g$ estrictamente decreciente;\quad $\operatorname{Rec}(g)=[0,\sqrt{3}]$

\end{enumerate}
\end{solbox}

\vspace{0.4cm}

% ── Sección II ────────────────────────────────────────────────
\begin{solbox}{II. Soluciones — Composición, Inyectividad y Función Inversa}
\begin{enumerate}[label=\color{azulFuerte}\bfseries 2.\arabic*., leftmargin=*, itemsep=0.6em, start=2]

    \item $\operatorname{Dom}(f\circ g)=[1,+\infty)$;\quad
          $(f\circ g)(5)=\sqrt{3}$;\quad
          Preimagen de $3$: $x=65$

    \item $\operatorname{Dom}(f\circ g)=(-\infty,0]\cup(2,+\infty)$;\quad
          $(f\circ g)(4)=\sqrt{2}$;\quad
          Preimagen de $1$: no existe

    \item $k=-\dfrac{3}{2}$

    \item $b=3$

    \item $g$ inyectiva en $[1,+\infty)$, no sobreyectiva\;
          $\bigl(\operatorname{Rec}(g)=[7,+\infty)\neq\mathbb{R}\bigr)$

    \item $f$ biyectiva;\quad
          $f^{-1}(y)=\dfrac{y-1+\sqrt{y^{2}+1}}{y}$,\quad
          $\operatorname{Dom}(f^{-1})=\mathbb{R}$,\;$\operatorname{Rec}(f^{-1})=\,]0,2[$

    \item $f$ biyectiva;\quad
          $f^{-1}(x)=2+\sqrt{\dfrac{x+5}{2}}$,\quad
          $\operatorname{Dom}(f^{-1})=[-5,+\infty)$,\;$\operatorname{Rec}(f^{-1})=[2,+\infty)$

\end{enumerate}
\end{solbox}

\vspace{0.4cm}

% ── Sección III — PEPs Antiguas ───────────────────────────────
\begin{solbox}{III. Soluciones — PEPs Antiguas}

{\color{azulFuerte}\textbf{PEP 2023}}
\begin{enumerate}[label=\color{azulFuerte}\bfseries\arabic*., leftmargin=*, itemsep=0.5em]
    \item[1.0)] $\operatorname{Dom}(f\circ g)=[-3,\,3]$

    \item[1.1)] $A=[2,+\infty)$,\; $B=(-\infty,1]$;\quad
                $f^{-1}(x)=2+\sqrt{1-x}$,\; $\operatorname{Dom}(f^{-1})=(-\infty,1]$
\end{enumerate}

\smallskip
{\color{azulFuerte}\textbf{PEP 2022}}
\begin{enumerate}[label=\color{azulFuerte}\bfseries\arabic*., leftmargin=*, itemsep=0.5em]
    \item[1.3)] $\operatorname{Dom}(f)=\bigl[-\tfrac{5}{2},\,-\tfrac{1}{2}\bigr]$;\quad
                $f$ decreciente;\quad $\operatorname{Rec}(f)=[-\sqrt{2},\,0]$
\end{enumerate}

\smallskip
{\color{azulFuerte}\textbf{PEP 2023 — Alternativa}}
\begin{enumerate}[label=\color{azulFuerte}\bfseries\arabic*., leftmargin=*, itemsep=0.5em]
    \item[1.4)] $A=[-4,+\infty)$,\; $B=[-2,+\infty)$;\quad
                $f^{-1}(x)=-4+\sqrt{x+2}$,\; $\operatorname{Dom}(f^{-1})=[-2,+\infty)$

    \item[1.5)] $\operatorname{Dom}(f)=\bigl(-\infty,\tfrac{3}{2}\bigr]\cup(4,+\infty)$;\;
                $\operatorname{Rec}(f)=[1,+\infty)$;\;
                Preimagen de $3$: $x=\dfrac{13}{2}$

    \item[1.6)] $f$ biyectiva;\quad
                $f^{-1}(y)=\dfrac{y}{1+|y|}$,\quad
                $\operatorname{Dom}(f^{-1})=\mathbb{R}$,\; $\operatorname{Rec}(f^{-1})=\,]{-1},1[$
\end{enumerate}

\smallskip
{\color{azulFuerte}\textbf{Más PEP}}
\begin{enumerate}[label=\color{azulFuerte}\bfseries\arabic*., leftmargin=*, itemsep=0.5em]
    \item[1.8a)] $k=2$;\; $f$ inyectiva y sobreyectiva;\quad
                 $f^{-1}(x)=2-\sqrt{x-1}$,\; $\operatorname{Dom}(f^{-1})=[2,10]$,\;
                 $\operatorname{Rec}(f^{-1})=[-1,1]$

    \item[1.8b)] $\operatorname{Dom}(f)=\bigl[-\tfrac{14}{3},\,\tfrac{2}{3}\bigr]$;\quad $f$ creciente

    \item[1.9)] $\operatorname{Dom}(f)=(-1,\,8]$,\; $\operatorname{Rec}(f)=[0,+\infty)$;\quad
                $(f\circ g)(x)=\sqrt{\dfrac{3}{\sqrt{x-1}}-1}$,\;
                $\operatorname{Dom}(f\circ g)=(1,\,10]$
\end{enumerate}

\end{solbox}

\vspace{0.8cm}
\begin{center}
    \small\textbf{Usach Premium} — ¡Nos vemos en la ayudantía! \\
    \textit{``Por y para estudiantes.''}
\end{center}

\end{soluciones}

\end{document}